\documentclass[10pt,a4paper]{amsart}
\usepackage[margin=19mm]{geometry}
\usepackage{amsmath,amssymb,mathtools}
\usepackage[scr=rsfs,cal=euler]{mathalfa}
\usepackage{xfrac}
\usepackage[unicode,hidelinks,bookmarks=false]{hyperref}
\newcommand{\ie}{i.e.\ }
\newcommand{\cf}{cf.\ }
\newcommand{\resp}{respectively\ }
\newcommand{\Subsection}{\subsection}
\providecommand{\nobreakdash}{\nobreak\hbox{-}}
\setlength{\emergencystretch}{2em}
\multlinegap0pt
\usepackage{amssymb}
\usepackage{mathtools}
\usepackage{enumerate}
\newtheorem{thm}{Theorem}
\newtheorem{prop}{Proposition}
\newcommand{\N}{\mathbb N}
\title{Sharp Fractional Riesz Estimates on the Hypercube}
\author{Yong Jiao, Sijie Luo, Dmitriy Zanin, Dejian Zhou}
\date{Source: arXiv:2609.09040v1 (CC BY 4.0)}
\begin{document}
\maketitle

\noindent\emph{Source and licence.} Sharp Fractional Riesz Estimates on the Hypercube. Version arXiv:2609.09040v1 (CC BY 4.0), distributed by arXiv under the Creative Commons Attribution 4.0 International licence, http://creativecommons.org/licenses/by/4.0/, as recorded in the arXiv licence metadata for this exact version and shown on the abstract page https://arxiv.org/abs/2609.09040v1. Full licence notice: \texttt{licenses/arxiv-2609.09040-CC-BY-4.0.txt}.\par\smallskip

\noindent\textit{Editorial scope.} Counted unit: the article's prop prop:tail-gradient-potential (source0.tex lines 1672--1683). The article's complete original proof of it is delivered with it (source0.tex lines 1685--1767, one \texttt{proof} environment of 83 lines). No mathematics is written, repaired or re-derived: the statement and the proof are the article's own lines, in the article's own notation, and no retained line is substituted or re-flowed.  Retained uncounted as the source-local context the proof needs: lines 296--302; lines 1664--1666.  Nothing was omitted from any retained span, so the package carries no omission record.  No retained line contains a citation, so the wrapper carries no bibliography and no bibliography entry.  No retained line contains a figure environment, an image-inclusion command or drawing code, so no image is delivered and the package declares no asset dependency.  Editorial formatting only, in addition: an AMS \texttt{amsart} preamble that restates the article's own declarations and macros used by the retained lines, namely the environments \texttt{thm}, \texttt{prop} on separate counters, together with the standard packages those lines need.  The retained environments print in this document as Theorem 1, Proposition 1 in order of appearance, whereas the article prints the same items with the numbering of its own full document.  The complete native source of the article as distributed in the arXiv e-print for this exact version is bundled unchanged as \texttt{sources/209771/source0.tex}.
\begin{thm}\label{thm:main}
For every $1<p\leq 2,$ for every
$n\in\N$ and every $f:\Omega_n\to\mathbb{C},$ we have
\[
\left\|\left(\sum_{j=1}^n|D_jf|^2\right)^{\frac12}\right\|_{L_p(\Omega_n)} \leq c_{{\rm abs}}(p-1)^{-2}\|\Delta^{\frac1p}f\|_{L_p(\Omega_n)}.
\]
\end{thm}

\begin{equation}\label{eq:tail-heat-smoothing}
 \|e^{-t\Delta}f\|_{L_{p}(\Omega_{n})}\le C_q\exp\!\left[-c_qd\min\left\{t,t^{1/\vartheta_p}\right\}\right]\|f\|_{L_{p}(\Omega_{n})},
\end{equation}

\begin{prop}
\label{prop:tail-gradient-potential}
Let $1<p<2$, let $p'=p/(p-1)$.  Then, for every $d\ge1$ and every
$h\in \mathcal{T}^{\geq d}(\Omega_{n})$,
\begin{equation}\label{eq:tail-gradient-potential}
 \|\nabla\Delta^{-1}h\|_{L_p(\Omega_n;\ell_2)}
 \lesssim_p
 d^{-\vartheta_p/p'}
 \|h\|_{L_p(\Omega_n)},
\end{equation}
where $\vartheta_q$ comes from \eqref{eq:tail-heat-smoothing}.
\end{prop}

\begin{proof}
By the sharp fractional Riesz estimate of Theorem \ref{thm:main}, we have
\begin{equation}
 \|\nabla u\|_{L_p(\Omega_n;\ell_2)}
 \lesssim_p
 \|\Delta^{1/p}u\|_{L_p(\Omega_n)},
 \qquad 1<p<2.
 \label{eq:critical-Riesz-tail-proof}
\end{equation}
Applying \eqref{eq:critical-Riesz-tail-proof} to
$u=\Delta^{-1}h$ gives
\begin{equation}\label{eq:gradient-potential-factor}
\begin{split}
\|\nabla\Delta^{-1}h\|_{L_p(\Omega_n;\ell_2)}
 &\leq \frac{c_{\mathrm{abs}}}{(p-1)^{2}}\|\Delta^{1/p}\Delta^{-1}h\|_{L_p(\Omega_n)}\\
 &=\frac{c_{\mathrm{abs}}}{(p-1)^{2}}\|\Delta^{-1/p'}h\|_{L_p(\Omega_n)},
\end{split}
\end{equation}
where we used $\frac{1}{p^{\prime}}=1-\frac{1}{p}$. Thus the problem is reduced exactly to a negative fractional-power estimate on the tail space.

We now put $\beta=\frac1{p^{\prime}}$. By Minkowski's inequality and
\eqref{eq:tail-heat-smoothing},
\begin{align}
 \|\Delta^{-\beta}h\|_{L_{p}(\Omega_{n})}
 &\le
 \frac{A_p}{\Gamma(\beta)}
 \|h\|_{L_{p}(\Omega_{n})}
 \int_0^\infty
 t^{\beta-1}
 \exp\!\left[
   -a_pd\min\{t,t^{1/\vartheta_p}\}
 \right]dt.
 \label{eq:negative-tail-integral}
\end{align}
We split the integral at $t=1$.  Since $\vartheta_p\le1$, for $0<t\le1$ we have $\min\{t,t^{1/\vartheta_p}\}=t^{1/\vartheta_p}$. Therefore, after the change of variables
$u=a_pd\,t^{1/\vartheta_p}$,
\begin{align}
 \int_0^1
 t^{\beta-1}e^{-a_pd t^{1/\vartheta_p}}\,dt
 &=
 \vartheta_p
 (a_pd)^{-\beta\vartheta_p}
 \int_0^{a_pd}
 u^{\beta\vartheta_p-1}e^{-u}\,du
 \notag\\
 &\le
 \vartheta_p
 \Gamma(\beta\vartheta_p)
 a_p^{-\beta\vartheta_p}
 d^{-\beta\vartheta_p}.
 \label{eq:small-time-tail-integral}
\end{align}
For $t\ge1$ the minimum equals $t$.  Since
$0<\beta<1$, we have $t^{\beta-1}\le1$ on $[1,\infty)$, and hence
\begin{align}
 \int_1^\infty
 t^{\beta-1}e^{-a_pd t}\,dt
 &\le
 \int_1^\infty e^{-a_pd t}\,dt
 =
 \frac{e^{-a_pd}}{a_pd}
 \lesssim_p
 d^{-\beta\vartheta_p}.
 \label{eq:large-time-tail-integral}
\end{align}
Combining
\eqref{eq:negative-tail-integral}--\eqref{eq:large-time-tail-integral}
gives
\begin{equation}
 \|\Delta^{-\beta}h\|_{L_{p}(\Omega_{n})}
 \lesssim_{p,\beta}
 d^{-\beta\vartheta_p}\|h\|_{L_{p}(\Omega_{n})}.
 \label{eq:negative-tail-beta}
\end{equation}
Recalling $\beta=1/p'$ and inserting \eqref{eq:negative-tail-beta} into
\eqref{eq:gradient-potential-factor} yields
\[
 \|\nabla\Delta^{-1}h\|_{L_p(\Omega_n;\ell_2)}
 \lesssim_p
 d^{-\vartheta_p/p'}\|h\|_{L_p(\Omega_n)},
\]
which proves \eqref{eq:tail-gradient-potential}.
\end{proof}

\end{document}
